% =============================================================================
% sec_related.tex -- Related Work  (ToG manuscript, item W6-related-work)
%
% Answers CoG 2026 reviews: R1 (value to the wider community, citing Schubert et al. 2016, Block et
% al. 2018, Kokkinakis et al. 2020, Pedrassoli Chitayat et al. 2024 (two papers) and Charleer et al.
% 2018; where the thresholds come from), R2 (why deep baselines, and why an equal-effort comparison),
% R3 (every algorithm motivated before it is used), R4 (fights without kills: pointer to the
% measurement).  Five paragraphs: (1) match and in-game win prediction; (2) encounter and team-fight
% detection; (3) audiences and broadcast; (4) deep learning versus gradient-boosted trees on tables;
% (5) win-probability models and calibration.
%
% Provenance conventions
%   * Every citation key exists in docs/references/definition_refs.bib, each marked [VERIFIED <date>]
%     there.  The comment after each sentence names the passage of the cited work that supports it and
%     the copy it was read on.  Copies other than the publisher's (arXiv, an authors' accepted
%     manuscript, a local PDF, an abstract served by the Semantic Scholar API) were used for content
%     only, never for bibliographic facts.  All passages were re-read on 2026-09-14 for this pass.
%       REFS = D:/LOL_Project/references (local PDFs, text extracted with pdftotext)
%       S2   = https://api.semanticscholar.org/graph/v1/paper/DOI:<doi>?fields=title,abstract
%   * Every number about this study carries its artefact and key.
%       FEAT = D:/LOL_Project/fusion_2615/features
%       REPO = C:/Users/todtj/PycharmProjects/LOL_teamfight
%       WT   = C:/Users/todtj/문서/LOL_Teamfight/worktrees/engagement-state-value
%   * No result of a run still in progress is stated or implied: the section makes no claim about
%     which learner is best, about scale classes, or about any calibration value.
%   * Game-version naming: no patch is named in the text.  Jalovaara's thesis names its data
%     "Patch 14.15" (Riot's pre-2025 displayed number); the text does not repeat it, so no mapping
%     between Riot's numbering schemes is relied on here (the mapping used elsewhere, displayed 25.N =
%     gameVersion 15.N, is documented in definition_refs.bib section 10).
%   * Cross-references use labels that exist on 2026-09-14 (sec:background-data, sec:definition,
%     sec:def-constants, sec:def-killless, sec:learners, sec:learners-input, sec:learners-effort,
%     subsec:future-work).  sec_limit.tex was not written when this file was drafted, so its label was rendered as
%     a visible pending-ref placeholder; item F4-consistency-fills (2026-09-14) replaced both with \ref{sec:limit}, whose
%     label exists (sec_limit.tex, \section{Where the Limit Comes From}).
%   * Paraphrase only; the text contains no direct quotation.
%   * Length: IEEEtran journal build on 2026-09-14 after the fact-check fixes (MiKTeX pdflatex + bibtex,
%     scratchpad copy with url, hyperref and cleveref as in main.tex, followed by sec_background,
%     sec_definition, sec_learners and sec_limitations so that cross-references resolve): no errors,
%     no undefined citations, 0 bibtex warnings, every \ref of this file resolved except the pending
%     sec:limit; about 1,100 words of prose, one page plus four lines while the two pending-ref
%     placeholders were rendered (both replaced by \ref{sec:limit} on 2026-09-14, item F4-consistency-fills).
%
% Triage of the previous draft (2026-09-14 05:41), kept where correct.  Changed in this pass:
%   - (2) added Maymin's own LoL engagement rule (damage-based, 5 s gap, experience-range gate), the
%     closest precedent for the rule anchors R and B, which the draft omitted;
%   - (2) the draft said all prior constants were "chosen by hand or computed from damage and position
%     streams"; Ke et al.'s radii are game ranges and the OpenDota rule reads only deaths, so the text
%     now says which rule needs which stream;
%   - (2) D is stated as in sec_definition.tex (sharing rate falls to 0.5), not "most pairs stop";
%   - (1) Bisberg and Ferrara are no longer called a pre-game model (their paper does not say so plainly);
%   - (4) the draft argued the data "differ in both respects" and used the 942 constant columns as the
%     second respect; constant columns are dropped before learning, so only size is used;
%   - (2) the timeline citation is attached to the timeline, not to the claim about its content, which
%     rests on sec_background.tex;
%   - (5) added Naeini et al.'s near-perfect-AUC, poorly calibrated example for "ranking is not
%     calibration";
%   - \ref{sec:limit} -> pending-ref placeholder (since restored to \ref{sec:limit}, item F4); prose shortened from about
%     1.1 pages to one page.
%
% Fact-check fixes (2026-09-14, second pass; all four findings re-verified and accepted):
%   - (2) "built from kills, as Ke et al.'s team fights are" misdescribed Ke et al., whose team fights
%     are position-built encounters filtered by a kill (Section VI-B item 7, VI-C); now "requires a kill,
%     as Ke et al.'s team fights do, but is built from kill events alone";
%   - (2) "these constants are set by hand" extended to OpenDota's 15 s and three deaths, for which the
%     source file gives no derivation either way; now the two are stated separately;
%   - (2) D is now "the distance at which the share of consecutive kill pairs within that gap that
%     involve a common champion falls to one half" (local, distance-binned rate, as sec_definition.tex);
%   - (4) "about 19 times that size, so none of them settles our case" was a non sequitur for Gorishniy
%     et al. and Shwartz-Ziv and Armon, whose benchmarks include larger datasets; the size point is now
%     tied to Grinsztajn et al. alone ("about 19 times that medium size") and the need to measure rests
%     on Gorishniy et al.'s finding ("with no family superior on every dataset").
%   - To stay near one page, two sentences of (2) were joined with a semicolon; no content was removed.
% =============================================================================

\section{Related Work}\label{sec:related}

% --- (1) Match-outcome and in-game win prediction ------------------------------------------
Outcome prediction for multiplayer online battle arena (MOBA) games such as League of Legends (LoL)
and Dota~2 has mostly targeted the match winner: from champion picks and the players'
past performance \cite{costa2021,hitar2023}, from a graph of a league's season \cite{bisberg2022}, or,
during the match, from per-minute statistics \cite{silva2018continuous} or live professional data
\cite{hodge2021win}.
% costa2021, S2 abstract (DOI 10.1109/CoG52621.2021.9619019): pre-game information for professional LoL
%   matches; historical performance the most accurate features.
% hitar2023, S2 abstract (DOI 10.1109/TG.2022.3153086): winner of professional LoL matches from pre-game
%   data only; built features include a player's skill with a champion and team synergies.
% bisberg2022, arXiv 2207.13191 abstract and Section IV-A: semi-supervised graph convolutional network
%   that learns the structure of a league over a season and classifies games by their neighbourhood
%   (nodes are team-games linked to the opponent and to the team's previous game).  Whether its inputs
%   are strictly pre-game is not stated plainly, so the text does not call it a pre-game model.
% silva2018continuous, publisher PDF p. 639 (abstract, Section I): recurrent networks, one input per
%   minute, 7621 competitive LoL matches.
% hodge2021win, S2 abstract (DOI 10.1109/TG.2019.2948469): live win prediction for professional DotA 2
%   matches, trialled at a major international tournament.
For in-game win probability in LoL, Maymin fitted a logistic regression on elapsed time and each
team's kills, turrets and large monsters, sampling one random minute per game because consecutive
minutes are strongly correlated \cite{maymin2021smart}; Jalovaara trained a deep network on 400,000
ranked matches from Riot's public application programming interface (API) \cite{jalovaara2024win}.
% maymin2021smart: open-access article (CC BY), copy https://d-nb.info/1367424143/34, Section 3.2
%   "In-Game Win Probability", pp. 17-18: features red and blue kills, towers and large monsters plus
%   minutes elapsed; logistic regression; one random minute per game because consecutive rows of the
%   same game are highly correlated.  Towers = turrets.
% jalovaara2024win: thesis PDF (url in the bib entry), abstract p. 3 (deep neural network for in-game
%   win probability) and Section 5.1: 400 000 Ranked matches collected with the LoL API; matches divided
%   into 5-minute intervals, one game state sampled uniformly per interval.
In Dota~2, deaths \cite{katona2019time} and team fights \cite{tot2021camera} have been predicted
seconds ahead, and team fights have fed match-winner prediction \cite{ke2022}, but none of these
studies predicts, before a fight in LoL begins, which team will win it.
% katona2019time, arXiv 1906.03939 abstract: death of a Dota 2 player within a five-second window.
% tot2021camera, REFS/Tot2021_TeamFightPrediction_CoG.pdf abstract and Section I: predicts and detects
%   team fights, shortly before their start, from player and camera positions (Dota 2).
% ke2022, REFS/Ke2022_TeamFightModels_CoG.pdf Sections VI-E and VII: team-fight features fed to recurrent
%   networks (GRU, LSTM) to classify the overall match outcome.
% "None of these": checked against the passages above for the nine works of this paragraph.  Schubert
%   et al. predict encounter outcomes in Dota 2; they are discussed in the next paragraph.

% --- (2) Encounter and team-fight detection ---------------------------------------------------
Predicting a fight's winner first needs a rule that delimits fights. Schubert et al.\ segmented 412
Dota~2 replays into \emph{encounters}, groups of opposing units within combat range (each unit's mean
damage distance plus one standard deviation) that continue while at least two opposing units stay
involved \cite{schubert2016encounter}. As 81.1\% of encounters had no kill, they predicted which team
killed more opponents only in decisive encounters, from what they describe as the initial state.
% schubert2016encounter: REFS/Schubert2016_EncounterDetection_SSAC.pdf (printed page numbers): Section
%   4.2 successor relation (timeout tau; at least two opposing units carried over), p. 7; combat range =
%   mean + 1 SD of each unit's damage distance, about 85% of damage, p. 9; 412 DOTA matches, p. 11;
%   18,744 of 23,110 encounters (81.1%) without kill events, 4,012 decisive encounters (one team killed
%   more opponents), p. 12; logistic regression predicting which team had more kills from number of
%   units, role distribution, gold and XP of the initial state, pp. 14-15.  No value of tau is given.
%   The Fig. 9 caption speaks of gold and XP "during the encounter" while the text calls them initial,
%   hence "what they describe as".
Ke et al.\ kept this structure, fixed the link radii at Dota~2's attack and healing ranges (700
and 400 units), and required a kill and at least two players on one side \cite{ke2022}. In LoL, Maymin
chained champion-on-champion damage into an engagement that ends after 5\,s without damage, a gap he
tied to the game's notion of combat, and added champions within experience range of it
\cite{maymin2021smart}.
% ke2022: REFS/Ke2022_TeamFightModels_CoG.pdf Section VI-B: combat link within the general attack range
%   (700 units), support link within the general healing range (400 units); team fight = encounter with
%   at least one kill and one side having at least two players of the same team.
% maymin2021smart, Section 3.2 "Basic and Advanced Stats", p. 19: an engagement is a continuous sequence
%   of champion-on-champion damage until 5 s pass with no damage; 5 s is consistent with the game's
%   in-combat concept (an item's speed bonus after 5 s out of combat); any champion within XP range of a
%   champion already in the engagement is included.  No numeric range is given in the paper.
These three rules need positions every tick or second, and Maymin's also damage, which he captured
from the game client \cite{maymin2021smart}; LoL's public Match-V5 timeline \cite{riotapi2024timeline} records
positions once a minute and no damage events (Section~\ref{sec:background-data}). The
OpenDota labels used by Tot et al.\ \cite{tot2021camera} need only deaths: in its public source, the
processor opens a fight 15\,s before a hero death, closes it after 15\,s without one, and keeps it if
it has at least three deaths \cite{opendota2018processteamfights}. Tot et al.\ found the labels' timing
inconsistent; the processor gives no derivation for its 15\,s or its three deaths, and our conference
version set its 18\,s and 4,000\,u by hand.
% "every tick or second": Schubert pp. 7-9 (stream of position updates per tick); Ke Section VI-C
%   (positions updated at each tick); Maymin p. 19 (optical tracking to know all champion locations
%   every second).  Maymin abstract and Section 1 (Introduction), p. 11: the publisher's API offers only
%   rudimentary, box-score-equivalent statistics; data captured with computer vision and client hooks.
% Timeline content: sec_background.tex subsection sec:background-data (frames once per minute; no
%   damage, ability-use or movement event; frames carry running damage totals, hence "no damage events").
% tot2021camera: REFS/Tot2021_TeamFightPrediction_CoG.pdf Section IV.A (team-fight labels of 1,457
%   professional matches from the OpenDota API) and Section VI (labels set too early or too late).
% opendota2018processteamfights: processors/processTeamfights.js at commit e2e4432, lines 10-57
%   (teamfightCooldown = 15; start = death time - 15; closes 15 after the last death; deaths >= 3; no
%   spatial condition, so the rule reads only deaths).  The repository cannot show which version
%   OpenDota's servers ran for Tot et al. (references_audit.md section 2a), hence "in its public
%   source".  Closing rule re-read on 2026-09-14 at
%   https://raw.githubusercontent.com/odota/core/e2e44328032e/processors/processTeamfights.js:
%   a fight closes when an event's time minus last_death >= teamfightCooldown (15).
%   "No derivation": the file declares the constant 15 and the filter deaths >= 3 with no comment or
%   derivation.  The text does not say they were set by hand, which the source does not show.
% 18 s and 4,000 u set by hand: sec_definition.tex opening paragraph (lines 14-15); CoG text
%   REPO/docs/CoG2026_Paper.md section 3.1.
Our engagement (Section~\ref{sec:definition}) requires a kill, as Ke et al.'s team fights do, but is
built from kill events alone, estimates its two boundaries from the corpus and re-estimates them for
every patch. The time gap lies at the valley between the two modes of the intervals between
consecutive kills, as session thresholds do for the gaps between user actions
\cite{halfaker2015session}, and the distance limit is the distance at which the share of consecutive
kill pairs within that gap that involve a common champion falls to one half. Its two remaining gates,
like Maymin's experience range, come from rules of the game (Section~\ref{sec:def-constants}), and
Section~\ref{sec:def-killless} measures what requiring a kill misses.
% Method statements: sec_definition.tex subsections "Kill episodes and the temporal boundary G" (KDE
%   valley between the two modes of the log kill-interval distribution), "The spatial boundary D"
%   (lines 118-125: among consecutive kill pairs no more than G apart, 5,460,008 pairs, the rate of
%   involving the same champion as killer, victim or assister is computed in distance bins, and D is
%   the distance at which that rate falls to 0.5; "share ... falls to one half" is this local,
%   per-distance rate, not a cumulative share), "Presence
%   gate, lead and anchor" (R = experience-sharing radius, B = assist window) and "Re-estimation per
%   patch" (G and D re-estimated per patch under a drift rule; sec_definition.tex motivates this with
%   chitayat2023beyond, which this section no longer cites for lack of space).
% "requires a kill, as Ke et al.'s team fights do, but is built from kill events alone": Ke et al.
%   (REFS/Ke2022_TeamFightModels_CoG.pdf) Section VI-B item 7 (team fight = encounter containing at
%   least one kill link) and Section VI-C (the detector reads a stream of player unit positions each
%   tick, links combat components into encounters, then keeps encounters with one or more kill links).
%   So the kill is a filter on position-built encounters; only the kill requirement is shared.  Our
%   engagement is built from kill events (sec_definition.tex, "Kill episodes and the temporal
%   boundary G" and "The spatial boundary D").
% halfaker2015session: REFS/Halfaker2015_UserSessionIdentification_WWW.pdf (arXiv copy) abstract and
%   Section 3.1: histogram of log inter-activity times inspected for a valley, two-component Gaussian
%   mixture, threshold where the within- and between-session components are equally likely.

% --- (3) Esports analytics for audiences and broadcast ---------------------------------------
Esports analytics includes communicating patterns found in data \cite{schubert2016encounter}, and
audience research builds such communication. At a Dota~2 tournament watched by 25 million people,
Echo, which turns exceptional player performances into broadcast graphics, affected the range and
quality of storytelling and raised audience engagement \cite{block2018narrative}. Spectators using
the Weavr companion app at two Dota~2 tournaments checked their reading of a match against its win
probability and doubted it when the two disagreed \cite{kokkinakis2020dax}; a streaming
extension that reached over 300,000 people showed win probability in place of the reliable score that
Dota~2 lacks \cite{pedrassoli2024passive}.
% schubert2016encounter p. 1 (Introduction): esports analytics defined to include communicating
%   patterns with visualization to assist decision-making.
% block2018narrative, S2 abstract (DOI 10.1145/3210825.3210833; ACM gold open access): Echo detects
%   extraordinary player performances in Dota 2 from live and historic data and turns them into
%   audience-facing graphics; tournament watched by 25 million; measurably affected range and quality of
%   storytelling; increased audience engagement.
% kokkinakis2020dax: authors' accepted manuscript, White Rose eprint 160015 (content only): abstract
%   (Weavr companion app, deployments at two major tournaments); Section 3.1.1 (Dota 2 has no score like
%   football); Section 3.3 (Win Probability view as a live score); Section 4.4.2 (Win Probability used
%   to check one's own intuition about the game); Section 5 (validity of win prediction questioned when
%   it conflicted with users' own perception).
% pedrassoli2024passive: authors' accepted manuscript, White Rose eprint 211999 (content only): abstract
%   (Dota 2 streaming extension reaching over 300,000 people) and Section 3.3 (Player Performance
%   panel: team win probabilities, given in lieu of a reliable score).
Spectator dashboards for LoL and Counter-Strike had to limit cognitive load and earn trust through
transparent design \cite{charleer2018dashboards}, and black-box win predictions leave commentators to
improvise explanations, which led to a model of the game state each team needs in order to win
\cite{pedrassoli2024wincondition}. For fights, this work points to two cues that viewers can check
against play: the probability, stated before an engagement begins, that each team wins it, and its
stakes, how far the match's win probability moves with its outcome
(Section~\ref{sec:limit}). Whether they help needs a user study
(Section~\ref{subsec:future-work}).
% charleer2018dashboards, S2 abstract (DOI 10.1145/3242671.3242680): dashboards for LoL and CS:GO
%   contribute to spectator insight and experience; in-game complexity and cognitive load must be
%   managed; trust through transparent design.
% pedrassoli2024wincondition, S2 abstract (DOI 10.1145/3677079): win predictions are mostly black-box
%   outputs that force commentators into ad-hoc interpretation; the win condition model predicts the
%   game state needed for each team to win.
% sec_limitations.tex subsection subsec:future-work, item "Audience-facing tools": no user study has
%   been run.  Stakes by pre-fight probability are the subject of sec_limit.tex (plan
%   REPO/docs/CLAUDE_TOG_PAPER_PLAN.md section 2, item 7.4); no value is stated here.

% --- (4) Deep learning and gradient-boosted trees on tabular data ----------------------------
For tables such as our pre-engagement input (Section~\ref{sec:learners-input}), it is open whether
neural networks or gradient-boosted decision trees, ensembles in which each tree corrects the errors
of those before it, work better. Papers proposing neural architectures for tables report gains on their
own benchmarks, among them TabNet, which selects features step by step
\cite{arik2021tabnet}, and SAINT, which relates features and rows to one another and outperforms
XGBoost, CatBoost and LightGBM on average \cite{somepalli2021saint}.
% arik2021tabnet, arXiv 1908.07442 abstract (preprint of the AAAI paper): sequential attention chooses
%   which features to reason from at each decision step; outperforms other neural network and decision
%   tree variants on non-performance-saturated tabular datasets.
% somepalli2021saint, arXiv 2106.01342 abstract: attention over both rows and columns; contrastive
%   self-supervised pre-training; outperforms XGBoost, CatBoost and LightGBM on average over a variety
%   of benchmark tasks.
Independent comparisons are less favourable. XGBoost beat recent deep models, even on the datasets of
the papers that proposed them, with much less tuning \cite{shwartzziv2022tabular}; after an extensive
hyper-parameter search on 45 datasets, tree-based models stayed ahead on medium-sized data of about
10,000 samples \cite{grinsztajn2022tree}; and under one training and tuning protocol, neither the best
deep models, among them the FT-Transformer, nor gradient-boosted trees were universally superior
\cite{gorishniy2021revisiting}. Our training patch alone holds 187,547 engagements, about 19 times that
medium size, and with no family superior on every dataset, the answer has to be measured on these
inputs. Section~\ref{sec:learners} therefore compares
a ladder of learners on identical inputs, each with a hyper-parameter search declared in advance
(Section~\ref{sec:learners-effort}), with LightGBM, which samples rows and bundles features to boost
quickly on large, wide data, as the tree learner \cite{ke2017lightgbm}.
% shwartzziv2022tabular, arXiv 2106.03253 abstract: XGBoost outperforms the deep models across the
%   datasets, including those used in the papers that proposed them; XGBoost requires much less tuning.
% grinsztajn2022tree, arXiv 2207.08815 abstract: 45 datasets; 20,000 compute hours of hyper-parameter
%   search per learner; tree-based models remain state-of-the-art on medium-sized data (~10K samples).
% gorishniy2021revisiting, arXiv 2106.11959 abstract: ResNet-like and Transformer-based (FT-Transformer)
%   models compared under the same training and tuning protocols; best DL models versus gradient boosted
%   decision trees: still no universally superior solution.
% 187,547: FEAT/model_comparison_v33_patch.json split.rows.train = 187547 (train patch 15.14), equal to
%   FEAT/model_comparison_input_audit.json truth.patches["15.14"].  187,547 / 10,000 = 18.75 -> "about 19".
%   Re-read 2026-09-14: split.rows = {train 187547, val 191184, test 153816}.
% Scope of the size point: "that medium size" refers back to the Grinsztajn et al. clause alone (the only
%   clause that names a size: medium-sized benchmark, ~10K samples, arXiv 2207.08815 abstract, re-read
%   2026-09-14).  "with no family superior on every dataset" restates Gorishniy et al.'s finding, cited
%   in the preceding sentence.  The size point does not separate our data from Gorishniy et al. (arXiv 2106.11959
%   PDF, dataset table row "#objects": 20,640 ... 500,000, 515,345, 581,012, 709,877, 1,200,192) or from
%   Shwartz-Ziv and Armon (arXiv 2106.03253 PDF, Section 3 and Table 1: 7,000 to 1,000,000 samples,
%   including Forest Cover Type 580k and Higgs Boson 800k).  Both PDFs re-read 2026-09-14.  For those two the text
%   instead uses Gorishniy et al.'s finding of no universally superior solution (abstract, re-read
%   2026-09-14): which family wins depends on the dataset, so it has to be measured on our inputs.
% Declared search: sec_learners.tex subsection sec:learners-effort (protocol only; its budget and all
%   results are \pending{learners} there).  No result is implied here.
% ke2017lightgbm, NeurIPS proceedings abstract page: GBDT implementations scale poorly when feature
%   dimension is high and data size is large; Gradient-based One-Side Sampling and Exclusive Feature
%   Bundling; up to over 20 times faster than conventional GBDT with almost the same accuracy.

% --- (5) Win-probability models and calibration ----------------------------------------------
A win probability measures value only if it is \emph{calibrated}: among situations given
probability $p$, the team should win about a fraction $p$ of them \cite{naeini2015obtaining}.
Maymin calls a kill smart and a death worthless by whether the team's estimated win
probability rose after it, and finds these track team performance far better than raw kills and deaths
\cite{maymin2021smart}; Jalovaara values items by win probability added, the change in win probability
that accompanies an action \cite{jalovaara2024win}.
% naeini2015obtaining, publisher PDF p. 2901 (Introduction): predictions of a binary outcome are well
%   calibrated if outcomes predicted with probability p occur about a fraction p of the time.
% maymin2021smart Section 3.2 "Worthless Deaths and Smart Kills", p. 18: a worthless death did not
%   increase the team's estimated win probability, a smart kill did; abstract: regular kills and deaths
%   do not nearly explain team performance as much as smart kills and worthless deaths.
% jalovaara2024win abstract (p. 3): win probability added = the change in win probability associated
%   with an action, a contextualised measure of value, applied to the evaluation of items.
Ranking well does not imply calibration: a classifier can separate two classes almost perfectly and
still be poorly calibrated, which post-processing by histogram binning, isotonic regression (a stepwise
mapping from scores to probabilities, fitted to the observed outcomes under the sole constraint that it
never decreases) or a
fitted sigmoid \cite{platt2000probabilities} is meant to correct \cite{naeini2015obtaining}. Guo et
al.\ found modern neural networks poorly calibrated and temperature scaling, a one-parameter variant
of Platt's sigmoid, effective on most datasets \cite{guo2017calibration}, and Kim et al.\ calibrated a
LoL winner predictor better than temperature scaling by modelling data uncertainty
\cite{kim2020confidence}.
% naeini2015obtaining p. 2903, Table 1(b) and its discussion: an SVM with AUC 1.00 and ECE 0.14 on a
%   simulated dataset; the learner discriminates very well but is poorly calibrated.  pp. 2901-2902:
%   histogram binning, Platt's sigmoid and isotonic regression as existing post-processing methods
%   (binning and isotonic regression non-parametric); ECE = sum over bins of the bin share times
%   |observed fraction - mean predicted probability|.
% Gloss of isotonic regression at its first use in the paper (item F3-fill-results, 2026-09-14; R3): a definition of the
%   method, not a claim of the cited works; the same wording is used in sec_limit.tex (sec:limit-value: "a stepwise mapping
%   from old to new probabilities that is constrained only never to decrease", cited there to zadrozny2002transforming).
%   No citation is added here.
% guo2017calibration, arXiv 1706.04599 abstract (preprint of the PMLR paper): modern neural networks,
%   unlike those from a decade ago, are poorly calibrated; on most datasets temperature scaling, a
%   single-parameter variant of Platt Scaling, is effective.
% platt2000probabilities: Platt's manuscript of the chapter dated 1999-03-26, copy at
%   https://home.cs.colorado.edu/~mozer/Teaching/syllabi/6622/papers/Platt1999.pdf (content only),
%   abstract: train an SVM, then train the parameters of an additional sigmoid that maps the SVM
%   outputs into probabilities.
% kim2020confidence, arXiv 2006.15521 abstract: data-uncertainty calibration for LoL winner prediction;
%   ECE 0.57% against 1.11% for temperature scaling.
Hence the AUC, which depends only on how predictions are ranked, is silent on calibration,
and a model calibrated on random match minutes need not be calibrated at engagement
cutoffs, which the kills that follow select; the value model of Section~\ref{sec:limit} is
therefore checked at those cutoffs.
% AUC is unchanged by any strictly increasing transformation of the scores (it depends only on their
%   ranking); stated without citation as a property of the measure.
% Check at engagement times exists: WT/scripts/validate_temporal_winprob_claude.py module docstring
%   (queries at engagement times are conditioned on a kill, so random-minute calibration does not
%   transfer by assumption; both populations scored with the same frozen models); results in
%   WT/outputs/temporal_winprob_v3_buckets/validation_a/results.json keys
%   engagement_all_available.pre and .last_kill_proxy (values cited in sec_intro.tex, not here).
